\section{Next Steps}
With more time, a non-linear model such as the random forest, gradient boosting, or interaction terms such as age $\times$ glucose, would be tested to address the linear decision boundary.
An L1 regularisation path would show which weak features could be removed, and the predicted probabilities could be calibrated so they can be read as a real stroke risk. Since the learning curve showed more rows alone are unlikely to help, new clinical features and an external stroke dataset could be used to improve the model and check whether it generalises.
